\section{Ejemplos de requisitos formalizados}
\label{sec:anexo-ejemplos}

A continuaci\'on se muestran dos ejemplos representativos de requisitos formalizados durante el trabajo.

\subsection{Ejemplo 1: Appwrite (funcional)}

\begin{lstlisting}[language=,caption={Ejemplo de requisito funcional --- Appwrite},label=lst:req-appwrite]
ID:
REQ-APPWRITE-10986

Nombre:
Bloqueo de creacion de usuarios sin verificacion

Tipo:
Funcional

Descripcion formal:
El sistema debera impedir la creacion de cuentas de usuario
cuando la direccion de correo electronico proporcionada no
haya sido verificada.

Fuente:
Release v6.0.0, PR #10986

Rationale:
Evitar la creacion de cuentas fraudulentas con correos no
verificados.

Prioridad:
Alta

Verificacion:
Intentar registrar un usuario con un correo no verificado
y comprobar que la API devuelve un error 403.

Estado:
Validado

Version:
1.0

Dependencias:

Modulo:
Auth / Users
\end{lstlisting}

\subsection{Ejemplo 2: n8n (restricci\'on)}

\begin{lstlisting}[language=,caption={Ejemplo de restricci\'on --- n8n},label=lst:req-n8n]
ID:
REQ-N8N-31371

Nombre:
Limite de nodos en workflows gratuitos

Tipo:
Restriccion

Descripcion formal:
El sistema debera limitar el numero maximo de nodos por
workflow a 10 para cuentas del plan gratuito.

Fuente:
Release v1.80.0, PR #31371

Rationale:
Control de recursos en el plan gratuito para garantizar
rendimiento del servicio.

Prioridad:
Alta

Verificacion:
Crear un workflow con 11 nodos desde una cuenta gratuita
y comprobar que el sistema rechaza la operacion.

Estado:
Validado

Version:
1.0

Dependencias:

Modulo:
Workflows / Plans
\end{lstlisting}
